\section{Experimental setup}
\label{sec:setup}

\subsection{Models and competence screen}
\label{sec:models}

We study two models most closely, the instruction-tuned Qwen3-8B \citep{yang2025qwen3} and Gemma 3 4B \citep{gemma2025gemma3}, loaded with int8 weights \citep{dettmers2022llmint8} and float16 (Qwen) or bfloat16 (Gemma) computation, using each model's chat template with thinking disabled for Qwen3; without a size, \emph{Qwen} and \emph{Gemma} refer to these two. Revisions and software are listed in Appendix~\ref{app:provenance}. A model enters the analysis only after passing a competence screen on a separate set of 24 histories, so that score sensitivity is studied while candidate selection is correct: in every construction, the correct current value must be strictly ranked first among the eight candidates on at least 99\% of the distinct screening prompts (ties count as errors), and in the superseded construction the mean margin of the current over the earlier value must be positive. Both models passed, with lowest construction-level accuracies of 287/288 (Qwen, live) and 575/576 (Gemma, entity mention). Seven further models were screened in the same way and the five that passed were run through Experiment 1 (Section~\ref{sec:additional-models}): Falcon3-7B-Instruct \citep{tii2024falcon3}, Granite-3.1-8B-Instruct \citep{granite2024granite3}, Mistral-7B-Instruct-v0.3 \citep{jiang2023mistral}, Gemma 3 12B and Qwen3-14B. OLMo-2-7B-Instruct \citep{olmo2025olmo2} and Qwen2.5-7B-Instruct \citep{qwen2025qwen25} failed, as did Phi-4-mini \citep{abouelenin2025phi4mini}; Llama-3.1-8B-Instruct was unavailable because access to its gated weights was declined.

\subsection{Analysis and prespecification}
\label{sec:inference}

All estimates are computed per history and then averaged over histories; histories are the unit of uncertainty. Intervals are percentile 95\% intervals from 2,000 bootstrap resamples of histories (seed 73021). Models are scored on the same histories and reported separately, without pooling.

The design, competence thresholds, and analysis of Experiment 1 were fixed in the project repository before its confirmation data were scored, and a dated amendment recorded before the results were examined specified the decision rule for the original hypothesis (Section~\ref{sec:relation}). The 24-history run of Experiment 2 was likewise specified before scoring. The 96-history run of Experiment 2 and the follow-up experiments were designed and their design documents hash-stamped on the experiment machine before any data were generated, but were not committed to version control. No experiment was externally preregistered. Order-stratified estimates, paired aligned-minus-reversed differences, per-attribute strata, and analyses restricted to correctly answered histories are secondary and exploratory.

\subsection{Experiments}
\label{sec:exp1}
\label{sec:marker-protocol}
\label{sec:followup-design}

\paragraph{Experiment 1: constructions crossed with order.} 96 histories: each of six entity pairs is crossed with the four target attributes and both assignments of names to the $x$ and $z$ roles, with two histories per cell; the other-attribute construction uses \lit{team} in half of them and \lit{project} in the other half. Each history yields $6\times2^5=192$ scored prompts (construction, query, edited entity, historical order, current order, baseline versus edit), 18,432 per model. Tokenizer-aware audits verified that each baseline/edit pair differs only at the intended one-token span. An earlier, simpler comparison that motivated this design, with a single unassigned control, is reported in Appendix~\ref{app:prelim}.

\paragraph{Experiment 2: the explicit marker.} In Experiment 1 the superseded and other-attribute constructions carry \lit{Previously} and the others do not. Experiment 2 crosses construction (superseded, entity mention) with marker (absent, present), where adding the marker changes nothing but the prefix \lit{Previously,} on each historical line (Table~\ref{tab:marker-templates}), with the other factors crossed as in Experiment 1. It uses 96 histories with one entity pair and the attribute \lit{badge}, so its $R$ values are compared within the experiment; an earlier 24-history run of the same design is a replication (Appendix~\ref{app:marker}). The prespecified quantities are the \emph{marker effect}, $R_{\text{present}}-R_{\text{absent}}$, within each construction, and the \emph{interaction}, the superseded minus the entity-mention marker effect.

\paragraph{Follow-ups.} Four follow-ups address questions left by Experiments 1 and 2, each with 96 histories, the same scoring and analysis, and a recorded prediction (designs in Appendix~\ref{app:followups}). \emph{Updated other attribute} adds a construction in which the other attribute is also updated, separating having been the queried attribute from having been contradicted by an update; the discount account predicts a negative update effect. \emph{Name--value distance} crosses construction with a near or far name--value distance, measured in tokens before data generation, with \lit{Previously} on every line; the prediction was near above far in both constructions. \emph{Update chains} give the queried attribute three to five successive values to find a depth at which answers fall below ceiling, and then estimate how $R$ grows with depth. An exploratory study of \emph{generated answers} uses OLMo-2-7B-Instruct, which fails the screen and answers below ceiling; its analogue of $R$ is \emph{donor following}, the share of edited prompts answered with the donor under the edited entity's query minus under the other entity's query.
